\section{Recurring Themes and Synthesis}\label{sec:recurring_themes}

Despite the diversity of cultural contexts, several recurring themes emerge in the treatment of dreams across civilizations, both in traditional and in industralized societies. These themes highlight the multifaceted roles that dreams have played in human life, encompassing practical, spiritual, and philosophical dimensions. Key recurring themes include:

\begin{itemize}
  \item \textbf{Healing and social regulation:} incubation rites, dietary and ritual prescriptions, and criteria for distinguishing “true” from “false” dreams \citep{lidonnici1995epidaurus,miller1994lateantiquity}.
  \item \textbf{Divination and decision-making:} from royal consultations in Mesopotamia to vernacular medieval manuals \citep{oppenheim1956dreams,cappozzo2018somniale}.
  \item \textbf{Training and ethics:} the contemplative discipline of dream yoga and the pedagogical role of dream-sharing in indigenous communities \citep{wangyal2001tibetan,tedlock1987anthrodreaming}.
  \item \textbf{Philosophy and identity:} metaphysical debates on reality versus illusion (China) and theological debates on divine versus demonic inspiration (Abrahamic traditions) \citep{campany2003dreamscape,miller1994lateantiquity}.
\end{itemize}